\documentclass[11pt]{amsart}
\usepackage{amssymb}
\usepackage{booktabs,array}
\usepackage[hidelinks]{hyperref}

\newtheorem{theorem}{Theorem}[section]
\newtheorem{lemma}[theorem]{Lemma}
\newtheorem{proposition}[theorem]{Proposition}
\theoremstyle{definition}
\newtheorem{remark}[theorem]{Remark}
\newtheorem{example}[theorem]{Example}
\numberwithin{equation}{section}

\newcommand{\N}{\mathbb{N}}
\newcommand{\Z}{\mathcal{Z}}
\newcommand{\Th}[1]{\Theta_{#1}}
\newcommand{\mo}{\bmod}

\title[Manabe's criterion for $b\le 2a$ and for $b\le 250$]{Necessity of Manabe's pure-periodicity criterion for three-move subtraction games with $b\le 2a$ or $b\le 250$}
\author{Aryan Poonacha}
\address{No affiliation. Contact: \url{https://github.com/Aryan-Poonacha/math-findings}}
\date{27 November 2026}
\subjclass[2020]{91A46; 11B50}
\keywords{subtraction game, Sprague--Grundy sequence, pure periodicity}

\begin{document}

\begin{abstract}
For the subtraction game $S=\{a,b,c\}$ with $2\le a<b<c$, $\gcd(a,b,c)=1$ and
$c\ne a+b$, Manabe (arXiv:2609.05358) gave three tests on $\rho=c\bmod(a+b)$ that imply
pure periodicity of the Sprague--Grundy sequence, and conjectured that they are also
necessary (his Conjecture~52). We prove this for every $b\le 2a$ and every $c$: three
explicit positions rule out every pure period. For fixed $a,b$, the P-position indicators at
$c+j$ and $2c+j$, for $-(a+b)\le j<2(a+b)$, depend only on $c\bmod(a+b)$ when $c\ge2(a+b)$.
This transfer, together with a finite witness search, proves the conjecture for every
$b\le250$. We also show that a game failing all three tests with
$c\ge 2b$ has no pure period $p\le c+b$.
\end{abstract}

\maketitle

\section{Introduction}

In the subtraction game with move set $S$, a position is a heap of $x\in\N=\{0,1,2,\dots\}$
tokens, a move removes $s\in S$ tokens with $s\le x$, and the player unable to move loses.
We consider $S=\{a,b,c\}$ with $0<a<b<c$. Its Sprague--Grundy sequence is eventually
periodic~\cite{Golomb}. Manabe~\cite{Manabe} gave a sufficient condition for it to be
periodic from position~$0$ and conjectured that the condition is necessary.

\subsection{The conjecture}
The following excerpts are from Manabe~\cite{Manabe}, version~2 (16 September 2026).
Section labels and editorial insertions appear in brackets; citations are omitted.
Here $[u,v)=\{u,\dots,v-1\}$ denotes an integer interval.

\begin{quote}\small
Its nim-value $G_S(x)$ is the least non-negative integer absent from the values of its
options, so $G_S(x)\in\{0,1,2,3\}$. The zero set $W_0(S)$ consists of the losing, or
P-positions. The nim-sequence is eventually periodic: $G_S(x+p)=G_S(x)$ for all
$x\geqslant q$. Taking $p\geqslant 1$ minimal and then $q\geqslant 0$ minimal defines the
period $p(S)$ and pre-period $q(S)$. Pure periodicity means $q=0$; we call the case $q>0$
pre-periodic. [\dots] We therefore focus on primitive rulesets ($\gcd(a,b,c)=1$) with
$2\leqslant a<b<c$ and $c\neq a+b$. [\S1]

[\S2] Let $B=S\setminus\{c\}=\{a,b\}$ and $\rho=c \bmod (a+b)$. [\dots] write
$\delta=b-a=\eta a+\varepsilon$, $\eta=\lfloor\delta/a\rfloor$, $0\leqslant\varepsilon<a$.
[\dots, Proposition 7] $\Z:=W_0(B)\cap[0,a+b)$ [\dots]

[\S3] [\dots] Translates of $\Z$ are read modulo $a+b$. [\dots] Let
$D=[a,b)\setminus W=[a,b)\setminus\Z$ be the base gaps below $b$. [Here $W=W_0(B)$.] [\dots]
\[
\Lambda=\begin{cases}\bigl[(\eta+1)a,\,(\eta+1)a+\varepsilon\bigr), & \eta \text{ even},\\
\bigl[\eta a+\varepsilon,\,(\eta+1)a\bigr), & \eta\text{ odd}.\end{cases}
\qquad\textup{(formula (4))}
\]
[\dots] \textsc{Definition 19.} Let $\rho\in\mathbb{Z}/(a+b)\mathbb{Z}$. We say that $\rho$ is
admissible for the period $P$, and write $\rho\in\Th{P}$, in the following three cases.
\begin{align*}
&(\mathrm{T0})\quad \rho\in\Th{a+b} :\Longleftrightarrow \Z\cap(\Z+\rho)=\varnothing;\\
&(\mathrm{T1})\quad \rho\in\Th{c+a} :\Longleftrightarrow \Z\cap[0,\delta)\cap\bigl(\Z+(\delta-\rho)\bigr)=\varnothing;\\
&(\mathrm{T2})\quad \rho\in\Th{c+b} :\Longleftrightarrow (\rho+\Lambda)\cap\Z=\varnothing \text{ and } (\rho+D)\cap\Lambda=\varnothing.
\end{align*}
All three are conditions on the angle alone; none of them involves $c$ beyond its residue.

\textsc{Conjecture 52.} Let $a\geqslant 2$, $\gcd(a,b,c)=1$ and $c\neq a+b$. Then $G_S$ is
purely periodic if and only if the angle is admissible for at least one candidate period,
that is, if and only if
\[\rho\in\Th{a+b}\cup\Th{c+a}\cup\Th{c+b},\]
and in that case the least period is the least admissible one,
\[p=\min\{P\in\{a+b,c+a,c+b\}:\rho\in\Th{P}\}.\]
Only necessity remains open; sufficiency and the least-period formula follow from Theorem 34.
\end{quote}

The sets $\Lambda$ and $D$ in (T2) are read modulo $a+b$ as well. Manabe's \S6 also states
Conjecture~56 (under the same hypotheses, the least period of a purely periodic $G_S$ is
one of $a+b$, $c+a$, $c+b$) and Conjecture~57 (if the least period is $c+a$, resp.\
$c+b$, then $\rho\in\Th{c+a}$, resp.\ $\rho\in\Th{c+b}$).

\subsection{Prior work}
We summarize results on the necessity direction and related period classifications.
\begin{itemize}
\item S.~Zhang~\cite[Prop.~3.1]{SZhang} determined the pre-period and period on the line
$b=2a$; with Manabe's Theorem~34 this proves Conjecture~52 there, as noted
in~\cite[\S6]{Manabe}. S.~Zhang~\cite[Prop.~4.1]{SZhang} also classified $b=a+1$, where
every angle is admissible. Three ultimately bipartite families~\cite[Thm.~6.3]{SZhang} and the
subfamily $b=a+2$ with $a$ odd are also known~\cite[\S6]{Manabe}.
\item Mikl\'os and Post~\cite{MP} proved that moves larger than a pure outcome period are
redundant~\cite[Prop.~2.9]{MP}, and gave a finite test for the outcome period
$b+c$~\cite[Cor.~2.12.1]{MP}.
\item Moriwaki (10 June 2026)~\cite[Thm.~3]{Moriwaki} treated real parameters with
$b\le 2a$. On his ``white regions'' he proved that none of $a+b$, $c+a$, $c+b$ is a pure
period, and conjectured that no pure period exists~\cite[Remark~5, \S5(1)]{Moriwaki}.
For integer parameters the white regions are exactly the non-admissible games with
$b\le 2a$ (Remark~\ref{rem:moriwaki}).
\item Manabe (v1 4 September, v2 16 September 2026)~\cite{Manabe} proved sufficiency and the
least-period formula (his Theorem~34), proved necessity for least periods
$p\le a+b$ when $c\ge 2(a+b)$ (his Proposition~55), and posed Conjectures~52, 56 and~57.
\item Later work gave exact nim-value words or period formulas for individual families:
$\{2,5,c\}$ with $c\ne7$ (Zeng, 8 September~\cite{Zeng}), $\{2,b,2b-6\}$ with odd $b\ge7$
(Jo, 9 September~\cite{Jo}), primitive $\{a,4a,c\}$ with $a\ge2$ and $c>4a$ (AI-generated
preprint, 15 October~\cite{A4AC}), $\{a,3a-1,c\}$ with $a\ge3$ and $c>(a-1)(4a-1)$ and
$\{6,17,119+23k\}$ with $k\ge0$ (Brink, 23 November~\cite{Brink1,Brink2}), and
$\{3,8,11k+7\}$ with $k\ge1$ (Cameron-Lyons, 23 November~\cite{CL}).
\item Levinstein (25 November 2026)~\cite{Levinstein} proved Conjecture~57 in a stronger
form, with a Lean formalization: for all $1\le a<b<c$ and $P\in\{a+b,c+a,c+b\}$, the
P-positions are periodic from~$0$ with period $P$ if and only if $\rho\in\Th{P}$. Hence,
under the hypotheses of Conjecture~52, it is equivalent to Conjecture~56. He also excluded periods
$p\in(c-a,c+a)$ \cite[Prop.~3.1]{Levinstein}, and reported a search of 100{,}946{,}926
triples with bounded $c$ without counterexample. H.~Zhang (26 November 2026)~\cite{HZhang}
gave a second proof of Conjecture~57 by a reflection argument, also with Lean.
\end{itemize}

\subsection{Results}
\begin{itemize}
\item \emph{Conjecture~52 holds whenever $b\le 2a$} (Theorem~\ref{thm:band}). If $b\le 2a$
and $\rho$ is not admissible, three explicit positions show that neither $G_S$ nor the set of
P-positions is periodic from~$0$, with any period. Manabe lists this region as open apart
from the cases above; the line $b=2a$ was known~\cite[Prop.~3.1]{SZhang}.
\item \emph{Conjecture~52 holds for every $b\le 250$ and every $c$} (Theorem~\ref{thm:250},
computer-assisted). The key step is a transfer lemma (Lemma~\ref{lem:transfer}): for fixed
$a,b$, the P-position indicators at $c+j$ and $2c+j$, $-(a+b)\le j<2(a+b)$, depend only on
$c\bmod(a+b)$ once $c\ge 2(a+b)$.
\item \emph{Towards Conjecture~56:} if $c\ge 2b$ and $\rho$ is not admissible, the
P-positions have no period $p\le c+b$ from~$0$ (Proposition~\ref{prop:cb}).
\end{itemize}
Theorems~\ref{thm:band} and~\ref{thm:250} use the obstruction in Lemma~\ref{lem:witness}.
Conjecture~52 remains open for
$b>\max(2a,250)$, and Conjecture~56 remains open.

\section{Definitions and known results}

Throughout, $S=\{a,b,c\}$ with $1\le a<b<c$; further hypotheses are stated where needed. Put
$m=a+b$, $\delta=b-a$ and $\rho=c\bmod m\in[0,m)$, and let $Y\colon\N\to\{0,1\}$ be the
indicator of the P-positions $W_0(S)$. Then
\begin{equation}\label{eq:rec}
Y(x)=1\iff Y(x-s)=0\text{ for every }s\in S\text{ with }s\le x,
\end{equation}
and $W_0(S)$ is the unique subset of $\N$ with this property~\cite[Lemma~4]{Manabe}. Let $z$
be the indicator of $W_0(B)$. Below $c$ the two games have the same moves, so
\begin{equation}\label{eq:below}
Y(x)=z(x)\qquad(0\le x<c).
\end{equation}
The base pattern is periodic: $z(x+m)=z(x)$ for all $x\ge0$, and
$\Z=\{x\in[0,m):z(x)=1\}$ satisfies $[0,a)\subseteq\Z\subseteq[0,b)$, with $\Z=[0,a)$ when
$b\le 2a$ \cite[Prop.~7]{Manabe}, \cite[Thms.~3.3--3.4]{MP}. We identify $\Z$ with a subset
of $\mathbb{Z}/m\mathbb{Z}$, so $z(x)=1$ if and only if $x\bmod m\in\Z$. Since
$b\equiv-a\pmod m$, the two moves of $B$ act on residues as $x\mapsto x-a$ and
$x\mapsto x+a$.

A function $f$ on $\N$ is \emph{periodic from $0$ with period $p\ge1$} if $f(x+p)=f(x)$ for
all $x\ge0$. Thus $G_S$ is purely periodic exactly when it is periodic from~$0$. Every period
of $G_S$ from~$0$ is a period of $Y$ from~$0$. We call $\rho$ (and $S$) \emph{admissible} if
$\rho\in\Th{a+b}\cup\Th{c+a}\cup\Th{c+b}$.

We use three known results.

\begin{theorem}[{Manabe~\cite[Thm.~34]{Manabe}}]\label{thm:34}
Let $a\ge2$, $\gcd(a,b,c)=1$ and $c\ne a+b$. If $\rho$ is admissible, then $G_S$ is purely
periodic and its least period is the least admissible candidate.
\end{theorem}

So Conjecture~52 holds for a given $S$ as soon as we know that $G_S$ is not purely periodic
when $\rho$ is not admissible.

\begin{theorem}[{Levinstein~\cite[Thm.~1.1]{Levinstein}; H.~Zhang~\cite{HZhang}}]\label{thm:57}
Let $P\in\{a+b,c+a,c+b\}$. If $Y$ is periodic from $0$ with period $P$, then
$\rho\in\Th{P}$.
\end{theorem}

\begin{lemma}[Move deletion; {\cite[Prop.~2.9]{MP}} and {\cite[Prop.~25]{Manabe}}]\label{lem:deletion}
If $Y$ is periodic from $0$ with a period $p\le c$, then $W_0(S)=W_0(B)$ and
$\rho\in\Th{a+b}$.
\end{lemma}

\begin{proof}
The first assertion is essentially~\cite[Prop.~2.9]{MP}; we include the short argument. Put
$W=W_0(S)$. By uniqueness~\cite[Lemma~4]{Manabe} for $B$, it suffices to show that $x\in W$ if and only if $x-s\notin W$ for
each $s\in\{a,b\}$ with $s\le x$. The forward implication holds by~\eqref{eq:rec}. For the
converse let $x\notin W$ and write $x=j+ip$ with $0\le j<p$. Then $j\notin W$, and the move
$c$ is not legal at $j<p\le c$, so $j-s\in W$ for some $s\in\{a,b\}$ with $s\le j$; by
periodicity $x-s\in W$. Finally $W_0(S)=W_0(B)$ implies
$\rho\in\Th{a+b}$~\cite[Prop.~25]{Manabe}.
\end{proof}

The following lemma gives the non-periodicity proofs of Theorems~\ref{thm:band}
and~\ref{thm:250}.

\begin{lemma}[Local obstruction]\label{lem:witness}
Write $S=\{s,t,u\}$ in any order and put $h=t+u$. If $Y$ is periodic from~$0$ (with any
period), then for every $x\ge h$ with $Y(x)=1$, exactly one of $Y(x-h)$ and $Y(x-h+s)$
equals~$1$. Consequently, if some $y\ge0$ satisfies
\begin{equation}\label{eq:witness}
Y(y)=0,\qquad Y(y+s)=0,\qquad Y(y+h)=1,
\end{equation}
then neither $Y$ nor $G_S$ is periodic from~$0$.
\end{lemma}

We call $(y,s)$ satisfying~\eqref{eq:witness} a \emph{witness}.

\begin{proof}
Let $p$ be a period of $Y$ from~$0$ and let $\overline W\subseteq\mathbb{Z}/p\mathbb{Z}$ be the
set of residues of P-positions, so $Y(x)=1$ if and only if $x\bmod p\in\overline W$. Applying
\eqref{eq:rec} to a representative $x\ge c$ of a residue $v$ gives
\begin{equation}\label{eq:circle}
v\in\overline W\iff v-a,\ v-b,\ v-c\notin\overline W\qquad(v\in\mathbb{Z}/p\mathbb{Z}).
\end{equation}
Let $v\in\overline W$. By~\eqref{eq:circle}, $v+t\notin\overline W$ and $v+u\notin\overline W$,
because both have the option $v$. The three options of $v+h$ are $v+u$, $v+t$ and $v+h-s$,
so by~\eqref{eq:circle} exactly one of $v+h$ and $v+h-s$ lies in $\overline W$ (these are
distinct residues: $\overline W\ni0$ is nonempty, and~\eqref{eq:circle} with $v\in\overline W$
shows $v-s\ne v$); call it
$\varphi(v)$. If $\varphi(v)=\varphi(v')$ with $v\ne v'$, then $v'=v\pm s$, and two elements
of $\overline W$ differ by the move $s$, contradicting~\eqref{eq:circle}. So $\varphi$ is an
injective map of the finite set $\overline W$ to itself, hence surjective. Thus every
$x\in\overline W$ has $x-h\in\overline W$ or $x-h+s\in\overline W$, and not both, since these
differ by $s$. For integers $x\ge h$ this is the first assertion. A witness violates it at
$x=y+h$. A period of $G_S$ from~$0$ is a period of $Y$ from~$0$.
\end{proof}

We have not found Lemma~\ref{lem:witness} in the literature; it is elementary and may be
known.

\section{Bases with \texorpdfstring{$b\le 2a$}{b <= 2a}}

In this section $a<b\le 2a$, so $1\le\delta\le a$ and $\Z=[0,a)$.

\begin{lemma}[{Manabe~\cite[Lemma~14, Prop.~26, \S6]{Manabe}}]\label{lem:angles}
If $a<b\le 2a$ then $\Th{a+b}=[a,b]$, $\Th{c+a}=\{0\}\cup[b,m)$ and $\Th{c+b}=[\delta,a]$.
Hence $\rho$ is admissible if and only if $\rho=0$ or $\rho\ge\delta$.
\end{lemma}

\begin{proof}
We include the short verification. On the circle $\mathbb{Z}/m\mathbb{Z}$, an arc
$[\sigma,\sigma+\ell)$ with $0\le\sigma<m$ misses $[0,\lambda)$, for $1\le\lambda,\ell<m$, if and only if
$\lambda\le\sigma\le m-\ell$. (T0) asks that $[\rho,\rho+a)$ miss $[0,a)$, that is,
$a\le\rho\le b$. In (T1), $\Z\cap[0,\delta)=[0,\delta)$, and with
$\sigma=(\delta-\rho)\bmod m$ the condition is $\delta\le\sigma\le b$; this holds for
$\rho=0$, fails for $1\le\rho\le\delta$ (then $\sigma<\delta$), and for $\delta<\rho<m$
(where $\sigma=\delta-\rho+m$) it means $\rho\ge b$. For (T2), $D=[a,b)$, and $\Lambda=[a,b)$
by formula~(4) (if $b<2a$ then $\eta=0$, $\varepsilon=\delta$; if $b=2a$ then $\eta=1$,
$\varepsilon=0$). So (T2) asks that the arc $\rho+[a,b)$ miss $[0,b)$; for $\rho<b$ this
means $\delta\le\rho\le a$, and for $\rho\ge b$ the arc starts at $\rho-b\in[0,a)$. The union
of the three sets is $\{0\}\cup[\delta,m)$.
\end{proof}

\begin{theorem}\label{thm:band}
Let $a<b\le 2a$, $c>b$ and $1\le\rho\le\delta-1$. Put $k=\lfloor c/m\rfloor$ and
$y=km-a-1$. Then $k\ge1$, and
\[
Y(y)=0,\qquad Y(y+c)=0,\qquad Y(y+m)=1 .
\]
Hence neither $Y$ nor $G_S$ is periodic from~$0$. In particular Conjecture~52 holds for every
$S$ with $b\le 2a$.
\end{theorem}

\begin{proof}
If $k=0$ then $c=\rho<\delta<b$, which is false; so $k\ge1$. Write $F(u)=Y(c+u)$ for
$u\ge-c$. We claim that for $0\le u<c$,
\begin{equation}\label{eq:H}
F(u)=1\iff u\bmod m\in H,\qquad H=[a,\,b-\rho)\ \cup\ [b,\,2a)\ \cup\ [m-\rho,\,m).
\end{equation}
Here $H\subseteq[a,m)$, and $[b,2a)$ is empty when $b=2a$. Let $0\le u<c$. The options of
$c+u$ are $u$, $c+u-a$ and $c+u-b$, and $Y(u)=z(u)=1$ exactly when $u\bmod m<a$.
Therefore
\begin{equation}\label{eq:Frec}
\begin{cases}
F(u)=0, & u\bmod m<a,\\
F(u)=1\iff F(u-a)=F(u-b)=0, & u\bmod m\ge a.
\end{cases}
\end{equation}

\emph{First block, $0\le u<m$.} For $u<a$, $F(u)=0$. For $a\le u<b$ we have
$u-a\in[0,\delta)\subseteq[0,a)$, so $F(u-a)=0$, while $c+u-b<c$ has value
$z(c+u-b)$ by~\eqref{eq:below}; note $c+u-b\equiv\rho+u-b$ and
$\rho+u-b\in[\rho-\delta,\rho)$. If $u\ge b-\rho$, this residue lies in
$[0,\rho)\subseteq[0,a)=\Z$, so the option is a P-position and $F(u)=0$. If $u<b-\rho$, the
residue is $\rho+u-b+m=\rho+u+a\in[2a,m)$, outside $\Z$, and $F(u)=1$. For $b\le u<m$ we have
$u-b\in[0,a)$, so $F(u-b)=0$ and $F(u)=1-F(u-a)$ with $u-a\in[\delta,b)$. If $u<2a$ then
$u-a<a$ and $F(u)=1$; if $u\ge2a$ then $u-a\in[a,b)$ and $F(u)=1$ if and only if
$u-a\ge b-\rho$, that is, $u\ge m-\rho$. This proves~\eqref{eq:H} for $u<m$; note $m<c$, as $k\ge1$ and $\rho\ge1$.

\emph{Later offsets, $m\le u<c$.} Now $u-a$ and $u-b$ are offsets in $[0,u)$, with residues
$x-a$ and $x+a$, where $x=u\bmod m$. By~\eqref{eq:Frec} and induction on $u$ it suffices to
check that for every $x\in[a,m)$,
\[
x\in H\iff (x-a)\bmod m\notin H\ \text{and}\ (x+a)\bmod m\notin H .
\]
If $a\le x<b$, then $x-a\in[0,a)$, and $x+a\in[2a,m)$ lies in $H$ exactly when $x\ge b-\rho$.
If $b\le x<2a$, then $x-a\in[\delta,a)$ and $(x+a)\bmod m=x-b\in[0,a)$, neither in $H$. If
$2a\le x<m$, then $x-a\in[a,b)$ lies in $H$ exactly when $x<m-\rho$, and
$(x+a)\bmod m=x-b\in[0,a)$. In each case the right side agrees with $x\in H$, which
proves~\eqref{eq:H}.

\emph{The witness.} We have $y=(k-1)m+b-1\ge0$ and $y<km\le c$. Since $y\bmod m=b-1\ge a$,
$Y(y)=z(y)=0$. Next $b-1\notin H$: $b-1\ge b-\rho$ as $\rho\ge1$, and $b-1<m-\rho$ as
$\rho<\delta\le a$. So $Y(y+c)=F(y)=0$. Finally $y+m=km+b-1=c+(b-\rho-1)$, and
$b-\rho-1\in[a,b-\rho)$ because $\rho\le\delta-1$; so $Y(y+m)=1$. This is a witness with
$s=c$ and $h=a+b=m$, and Lemma~\ref{lem:witness} applies.

For the last statement, let $S$ satisfy the hypotheses of Conjecture~52 with $b\le2a$. If
$\rho$ is admissible, Theorem~\ref{thm:34} applies. Otherwise $1\le\rho\le\delta-1$ by
Lemma~\ref{lem:angles}, and $G_S$ is not purely periodic.
\end{proof}

The first part of Theorem~\ref{thm:band} does not use primitivity, and its hypothesis forces
$a\ge\delta\ge2$ and $c\ne a+b$.

\begin{example}
For $S=\{5,8,28\}$: $m=13$, $\delta=3$, $\rho=2$, $k=2$, $y=20$. The positions $20$ and $48$
are N-positions and $33$ is a P-position, so no pure period exists.
\end{example}

\begin{remark}[Moriwaki's white regions]\label{rem:moriwaki}
In~\cite{Moriwaki} the largest move is normalized to~$1$. The $n$-th white region is
$na+nb<1<(n-1)a+(n+1)b$ \cite[Thm.~3]{Moriwaki}. At the point $(a/c,b/c)$, with integers
$a<b\le2a$ and $c>b$, this reads $nm<c<nm+\delta$, that is, $1\le\rho\le\delta-1$ and $n=k$.
By \cite[Prop.~2]{Moriwaki} and rescaling, his nim-value function at this point is
$f(x)=G_S(\lfloor cx\rfloor)$ for real $x\ge0$. Suppose $f(x+P)=f(x)$ for all $x\ge0$, and
put $N=\lfloor cP\rfloor$. If $cP\notin\mathbb{Z}$, comparing $x=n/c$ with $x$ slightly below
$(n+1)/c$ gives $G_S(n+N)=G_S(n)=G_S(n+N+1)$ for all $n\ge0$, so $G_S$ would be eventually
constant, which is impossible: a position whose options all have value $v$ does not
have value~$v$. So $cP$ is an integer and a pure period of $G_S$. Theorem~\ref{thm:band}
proves the non-pure-periodicity assertion of Moriwaki's Remark~5~\cite{Moriwaki} at rational
points.
\end{remark}

\section{Bases with \texorpdfstring{$b\le 250$}{b <= 250}}

\begin{lemma}[Transfer]\label{lem:transfer}
Let $1\le a<b$, $m=a+b$ and $0\le r<m$. For $k\ge2$ let $c_k=km+r$ and let $Y_k$ be the
P-position indicator of $\{a,b,c_k\}$. Then for every $j$ with $-m\le j<2m$, the values
$Y_k(c_k+j)$ and $Y_k(2c_k+j)$ do not depend on $k\ge2$.
\end{lemma}

\begin{proof}
Fix $k\ge2$, write $c=c_k$, $Y=Y_k$ and $F(u)=Y(c+u)$. We show that each value in question is
determined by $a,b,r,j$ through steps that do not involve $k$.

(a) For $-m\le u<0$, $F(u)=z(c+u)$ by~\eqref{eq:below}, and $c+u\equiv r+u\pmod m$.

(b) For $0\le u<c$ with $u\bmod m\in\Z$, $F(u)=0$, since the option $u$ of $c+u$ has
$Y(u)=z(u)=1$.

(c) For $0\le u<m$, the options of $c+u$ are $u$, with $Y(u)=z(u)$, and $c+(u-a)$, $c+(u-b)$
with $u-a,u-b\in[-m,u)$. By (a) and induction on $u$, $F(u)$ is determined by $a,b,r,u$.

(d) Let $N=(\mathbb{Z}/m\mathbb{Z})\setminus\Z$, with residues represented in $[0,m)$.
(i) Every $x\in N$ has $x-a\in\Z$ or $x+a\in\Z$: apply the recurrence of $B$ at
$x+m\ge b$, whose options $x+m-a$ and $x+m-b$ have residues $x-a$ and $x+a$.
(ii) If $e\in N$ and $e+a\in N$, then $\delta\le e<b$: for $e\ge b$,
$(e+a)\bmod m=e-b\in[0,a)\subseteq\Z$; for $e<\delta$, the position $e+a<b$ has only the move
$a$ in $B$, so $z(e+a)=1-z(e)=1$. Let $E=\{e\in N:e+a\in N\}$. By (i), each $e\in E$ has
$e-a\in\Z$ and $e+2a\in\Z$. Hence the pairs $\{e,e+a\}$, $e\in E$, are pairwise disjoint (if
$e'=e+a$ for $e,e'\in E$, then $e+a$ would have both neighbours $e$, $e+2a$ in $N$), and $N$
is the disjoint union of $E$, $E+a$ and $I=\{x\in N:x-a\in\Z,\ x+a\in\Z\}$. All additions
and subtractions in the definitions of $E$ and $I$ are modulo $m$; Manabe identifies $E$
with $\Lambda$~\cite[Lemma~22]{Manabe}. By (ii),
$E\subseteq[\delta,b)$ and $E+a\subseteq[b,m)$ as integers, so these residues lie in the
first block $[0,m)$.

(e) For $e\in E$, $F(e+a)=1-F(e)$. Indeed, the options of $c+e+a$ are $e+a$, an N-position
because $e+a<m\le c$ and $e+a\notin\Z$; $c+e$; and $c+e+a-b=c+(e-\delta)$ with
$e-\delta\in[0,a)\subseteq\Z$, an N-position by~(b).

(f) Let $H=I\cup\{x\in E\cup(E+a):F(x)=1\}$. By (c) and (e), $H$ is determined by $a,b,r$
and contains exactly one element of each pair $\{e,e+a\}$. We claim $F(u)=1$ if and only if
$u\bmod m\in H$, for $m\le u<c$. The options of $c+u$ are $u$, $c+(u-a)$ and $c+(u-b)$.
Write $x=u\bmod m$. The offsets $u-a$ and $u-b$ have residues $x-a$ and $x+a$, respectively,
and satisfy $0\le u-b<u-a<u$. We argue by
induction on $u$. If $x\in\Z$, then $F(u)=0$ by~(b). Otherwise the option $u$ is an
N-position, so $F(u)=1$ if and only if $F(u-a)=F(u-b)=0$. If $x\in I$, both $F(u-a)$ and
$F(u-b)$ vanish by~(b), so $F(u)=1$. If $x=e+a$ with $e\in E$, then $F(u-b)=0$ by~(b)
since $x+a=e+2a\in\Z$, and $u-a$ has residue $e$ and $u-a\ge m$ (as $u\ge m+e+a$); so
$F(u)=1-F(u-a)=1-[e\in H]=[e+a\in H]$. If $x=e\in E$, then $F(u-a)=0$ by~(b) since
$e-a\in\Z$, so $F(u)=1-F(u-b)$, where $u-b\equiv e+a$. If $u-b\ge m$, induction gives
$F(u)=1-[e+a\in H]=[e\in H]$. Otherwise $u=e+m$ and $u-b=e+a$, and by~(e)
$F(u)=1-F(e+a)=F(e)=[e\in H]$.

(g) Since $k\ge2$, we have $c\ge 2m$ and $c-m\ge m$. For $-m\le j<0$, $Y(c+j)$ is given
by~(a); for $0\le j<m$ by~(c); for $m\le j<2m\le c$ by~(f). Next,
$Y(2c+j)=F(c+j)$, and for $-m\le j<0$ we have $c+j\in[m,c)$, so (f) gives
$F(c+j)=[(r+j)\bmod m\in H]$. For $0\le j<2m$ the options of $2c+j$ are $c+j$, already
determined, and $2c+(j-a)$, $2c+(j-b)$ with $j-a,j-b\in[-m,j)$; induction on $j$ completes
the proof.
\end{proof}

\begin{theorem}\label{thm:250}
Conjecture~52 holds for every $S=\{a,b,c\}$ with $2\le a<b\le 250$, $c>b$,
$\gcd(a,b,c)=1$ and $c\ne a+b$.
\end{theorem}

\begin{proof}
This proof is computer-assisted (Section~\ref{sec:cert}). By Theorem~\ref{thm:34} we must
show that $G_S$ is not purely periodic when $\rho$ is not admissible. Note
$\gcd(a,b,c)=\gcd(a,b,\rho)$, since $\gcd(a,b)\mid m$.

\emph{Games with $c\ge 2m$.} Write $c=km+r$ with $k\ge2$ and $r=\rho$. For each base with
$2\le a<b\le250$ and each $r\in[0,m)$ with $\gcd(a,b,r)=1$, the program computes the angle
tests of Definition~19. Of these classes $(a,b,r)$, 3{,}902{,}781 are admissible. For each
of the other 2{,}568{,}108 it computes $Y$ for $c_0=2m+r$ from~\eqref{eq:rec} and finds a
witness $(y,s)$ with $y=c_0+j$, $-m\le j<m$, and $s\in\{a,b,c_0\}$. Every position in the
witness is of the form $c_0+i$ or $2c_0+i$ with $-m\le i<2m$: if $s=c_0$ (so $h=m$), these
are $c_0+j$, $2c_0+j$ and $c_0+(j+m)$; if $s\in\{a,b\}$ (so $h=m-s+c_0$), they are $c_0+j$,
$c_0+(j+s)$ and $2c_0+(j+m-s)$. By Lemma~\ref{lem:transfer}, for every $c=km+r$ with $k\ge2$
the same $j$, with the move $s$ if $s\in\{a,b\}$ and with $c$ in place of $c_0$ otherwise,
is a witness for $\{a,b,c\}$. Lemma~\ref{lem:witness} excludes pure periodicity.

\emph{Games with $b<c<2m$.} Among the primitive triples with $2\le a<b\le250$, $b<c<2m$ and
$c\ne m$, the program finds 5{,}579{,}832 admissible ones and, for each of the other
3{,}020{,}764, a witness computed directly from~\eqref{eq:rec}. Lemma~\ref{lem:witness}
applies again.
\end{proof}

Theorem~\ref{thm:band} covers $b\le 2a$ by proof, so the computation is needed only for
$2a<b\le250$; the program checks all bases anyway.

\section{Periods up to \texorpdfstring{$c+b$}{c+b}}

\begin{proposition}\label{prop:cb}
Let $1\le a<b<c$ with $c\ge 2b$, and suppose $\rho$ is non-admissible. Then $Y$ has no
period $p\le c+b$ from~$0$. The same exclusion holds for $G_S$.
\end{proposition}

\begin{proof}
Note $c\ge2b>a+b$. If $b\le2a$, Lemma~\ref{lem:angles} and Theorem~\ref{thm:band} apply, so let $b>2a$ and suppose
$Y$ has period $p\le c+b$ from~$0$. If $p\le c$, Lemma~\ref{lem:deletion} gives
$\rho\in\Th{a+b}$, a contradiction. So $c<p\le c+b$; put $d=p-c\in[1,b]$.

No two P-positions differ by $d$: if $Y(x)=Y(x+d)=1$, then $Y(x+p)=1$ and
$(x+p)-(x+d)=c$, contradicting~\eqref{eq:rec}. Let $z_0\in\Z$. Then $z_0<b$ and
$z_0+d<2b\le c$, so $Y(z_0)=z(z_0)=1$ and $Y(z_0+d)=z(z_0+d)$ by~\eqref{eq:below}. Hence
$(z_0+d)\bmod m\notin\Z$ for every $z_0\in\Z$, that is, $\Z\cap(\Z+d)=\varnothing$ and
$d\bmod m\in\Th{a+b}$ by (T0).

If $b=ka$ with $k$ odd, then $\eta=k-1$ is even and $\varepsilon=0$, so $\Lambda=\varnothing$
by formula~(4) and (T2) holds for every angle, contradicting non-admissibility. Otherwise
$\Th{a+b}=\{a,b\}$ by~\cite[Lemma~14]{Manabe}, where $\Th{a+b}$ is denoted $I$. Since
$1\le d\le b<m$, we get $d\in\{a,b\}$, so $p\in\{c+a,c+b\}$, and Theorem~\ref{thm:57} gives
$\rho\in\Th{p}$, again a contradiction.
\end{proof}

\begin{remark}
Manabe's Proposition~55 excludes least periods $p\le a+b$ when $c\ge2(a+b)$, and
Levinstein~\cite[Prop.~3.1]{Levinstein} excludes $p\in(c-a,c+a)$ for all $c$.
Under the hypotheses of Conjecture~52, any counterexample with $c\ge2b$ has least period
larger than $c+b$. The proof uses the complementary-difference condition (see
Manabe~\cite[\S3]{Manabe} and Levinstein~\cite[Lemma~2.3]{Levinstein}) on the unchanged
two-move prefix to force $(p-c)\bmod m\in\Th{a+b}$.
\end{remark}

\section{Computational certificates}\label{sec:cert}

All computations are reproduced by \texttt{verify.py}, a Python~3 program that uses only the
standard library and is available at
\url{https://github.com/Aryan-Poonacha/math-findings} (directory
\texttt{subtraction-games}). It computes P-positions from~\eqref{eq:rec}, in blocks
of $a$ consecutive positions with exact integer bit operations, and compares the bit-parallel
generator with mex on the first row's finite range. It computes $\Z$ from the two-move recurrence, $\Lambda$ from formula~(4), and the
sets $\Th{P}$ literally from Definition~19. It reads no stored results; any mismatch stops it
with an error. Table~\ref{tab:checks} lists the checks.

\begin{table}[!ht]
\small
\caption{Checks in \texttt{verify.py}; ranges in the default mode and with \texttt{--full}.}
\label{tab:checks}
\begin{tabular}{@{}>{\raggedright\arraybackslash}p{0.43\textwidth}>{\raggedright\arraybackslash}p{0.24\textwidth}>{\raggedright\arraybackslash}p{0.24\textwidth}@{}}
\toprule
Check & default & \texttt{--full}\\
\midrule
Bit-parallel recurrence equals the zero set of the mex rule & $1\le a\le29$,
$a<b<c\le40$, $0\le x<3c+7$ & $1\le a\le29$, $a<b<c\le40$, $0\le x<3c+7$\\
Manabe's box (primitive non-additive triples, $2\le a\le8$, $b\le40$, $c\le100$): exact
nim cycles; 7{,}557 admissible games pure with the predicted least period, 7{,}828
non-admissible games not pure, each with a witness; no witness in any pure game, with
witness search $0\le y<5c-m$ & all & all\\
Lemma~\ref{lem:angles}; formula~\eqref{eq:H} and the witness of Theorem~\ref{thm:band}
for $k=1,2,3$ & $b\le100$; $a\le25$ & $b\le250$; $a\le40$\\
Lemma~\ref{lem:transfer}: identical windows for $k=2,3,4,5$ & $b\le30$, all $r$, $a\ge1$
& $b\le60$\\
Theorem~\ref{thm:250}, $c\ge2m$: class witnesses at $c_0=2m+r$ & $b\le100$, plus direct
recheck at $k=3,4$ & $b\le250$, plus direct recheck at $k=3$\\
Theorem~\ref{thm:250}, $b<c<2m$: direct witnesses & $b\le100$ & $b\le250$\\
Angle facts used in Proposition~\ref{prop:cb} & $2\le a<b\le100$ & $2\le a<b\le250$\\
\bottomrule
\end{tabular}
\end{table}

The full run asserts the four counts in the proof of Theorem~\ref{thm:250}. The direct
rechecks evaluate the transferred witnesses at $c=3m+r$ (and, in the default mode, at
$c=4m+r$) without using
Lemma~\ref{lem:transfer}. With 16 processes and Python~3.12 the default mode takes about
10 seconds and the full mode about 70 seconds (8 minutes of processor time); the full log is
\texttt{logs/verify-full.log}.

The same four counts were obtained independently by a C++ program written from the
definitions during the audit described in Section~\ref{sec:ai}.

\section{Open questions}

\begin{enumerate}
\item Conjecture~52 for $b>\max(2a,250)$. For a fixed base, Lemma~\ref{lem:transfer}
reduces it to finitely many games whenever each non-admissible class has a witness in the
window of the lemma, as happens for $b\le250$. One approach for $b>2a$ is to describe the
first affected block and the set $H$ of Lemma~\ref{lem:transfer} explicitly.
\item Under the hypotheses of Conjecture~52, does every game that fails (T0)--(T2) have a
witness in the sense of Lemma~\ref{lem:witness}? The certificate finds one for each
non-admissible game in its range.
\item Conjecture~56 when $b<c<2b$ or $p>c+b$ (compare Proposition~\ref{prop:cb}).
\item Moriwaki's conjecture~\cite[Remark~5]{Moriwaki} at irrational points of the white
regions, and his assertion there that each white region splits into infinitely many parts.
\end{enumerate}

\section{Use of AI and verification}\label{sec:ai}

The author used AI systems for mathematical exploration, programming, drafting and checking.
A separate model-assisted audit compared the proofs with Manabe's preprint and recomputed the
certificate counts with an independent implementation. The author takes responsibility for
the results and the text.

The proofs of Lemmas~\ref{lem:witness} and~\ref{lem:transfer}, Theorem~\ref{thm:band} and
Proposition~\ref{prop:cb} are not computer-assisted; \texttt{verify.py} checks them only on
the finite ranges of Table~\ref{tab:checks} (\texttt{python3 -B verify.py}). The proof of
Theorem~\ref{thm:250} is computer-assisted: its two finite parts are reproduced by
\texttt{python3 -B verify.py --full}, and the default mode reproduces them for $b\le100$.

\begin{thebibliography}{99}
\raggedright

\bibitem{Brink1}
A.~Brink, \emph{Exact preperiods and periods for subtraction games with moves
$\{a,3a-1,c\}$}, Zenodo, 23 November 2026,
\href{https://doi.org/10.5281/zenodo.24485895}{doi:10.5281/zenodo.24485895}.

\bibitem{Brink2}
A.~Brink, \emph{Exact periods and insertion rules for the subtraction games
$\{6,17,119+23k\}$}, Zenodo, 23 November 2026,
\href{https://doi.org/10.5281/zenodo.24485055}{doi:10.5281/zenodo.24485055}.

\bibitem{CL}
Cameron-Lyons (GitHub user), \emph{formal-math-lab}, directory \texttt{manabe-residue-seven}
(Lean), 23 November 2026, \url{https://github.com/Cameron-Lyons/formal-math-lab}, accessed
27 November 2026.

\bibitem{A4AC}
\emph{A complete word classification for subtraction games with moves $a,4a,c$},
AI-generated preprint, version 1.1.1, 15 October 2026,
\url{https://github.com/TommyMaHaoKun/subtraction-games-a4ac/releases/tag/v1.1.1}, accessed
27 November 2026.

\bibitem{Golomb}
S.~W.~Golomb, \emph{A mathematical investigation of games of ``take-away''}, J. Combin.
Theory \textbf{1} (1966), 443--458.

\bibitem{Jo}
M.~Jo, \emph{Exact Grundy sequences for three-move subtraction games $\{2,b,2b-6\}$ with odd
$b$}, Zenodo, 9 September 2026,
\href{https://doi.org/10.5281/zenodo.22669492}{doi:10.5281/zenodo.22669492}.

\bibitem{Levinstein}
S.~Levinstein, \emph{Admissible angles for purely periodic three-move subtraction games},
Zenodo, 25 November 2026,
\href{https://doi.org/10.5281/zenodo.24536124}{doi:10.5281/zenodo.24536124}; Lean code at
\url{https://github.com/s66lem/conjectures}, accessed 27 November 2026.

\bibitem{Manabe}
H.~Manabe, \emph{Purely periodic three-move subtraction games}, preprint,
\href{https://arxiv.org/abs/2609.05358v2}{arXiv:2609.05358v2}, 2026.

\bibitem{MP}
I.~Mikl\'os and L.~Post, \emph{Superpolynomial period lengths of the winning positions in the
subtraction game}, Internat. J. Game Theory \textbf{53} (2024), 1275--1313,
\href{https://doi.org/10.1007/s00182-024-00911-5}{doi:10.1007/s00182-024-00911-5}; result
numbers follow \href{https://arxiv.org/abs/2312.02426v1}{arXiv:2312.02426v1}.

\bibitem{Moriwaki}
Y.~Moriwaki, \emph{Subtraction Nim with continuous parameters}, preprint,
\href{https://arxiv.org/abs/2606.11658v1}{arXiv:2606.11658v1}, 2026.

\bibitem{SZhang}
S.~Zhang, \emph{On the linearity of the periods of subtraction games}, Theoret. Comput. Sci.
\textbf{985} (2024), 114350,
\href{https://doi.org/10.1016/j.tcs.2023.114350}{doi:10.1016/j.tcs.2023.114350}; result
numbers follow the
\href{https://zhangshenxing.github.io/publications/Zhang2024tcs\%20On\%20the\%20linearity\%20of\%20the\%20periods\%20of\%20subtraction\%20games.pdf}{author's
copy}.

\bibitem{HZhang}
H.~Zhang, \emph{Pair-sum periods in three-move subtraction games}, manuscript with Lean
formalization, 26 November 2026,
\url{https://github.com/Z-bitlance/three-move-subtraction-periods}, accessed 27 November 2026.

\bibitem{Zeng}
Z.~Zeng, \emph{Pure periodicity for the three-move subtraction games $\{2,5,c\}$}, Zenodo,
8 September 2026,
\href{https://doi.org/10.5281/zenodo.22661435}{doi:10.5281/zenodo.22661435}.

\end{thebibliography}

\end{document}
